% !TeX root = ../抽象代数笔记.tex
% 由主文档通过 \input 载入；请编译主文档。
\setcounter{chapter}{9} % 保留原书章号。
\chapter{环、整环、域与环同态}\label{ch:rings}

\begin{description}
  \item[学习目标] 区分环、整环与域，将群同态的思想推广到同时保持两种运算的映射。
  \item[先修知识] 第\ref{ch:quotient-groups}章。
\end{description}

\section{环的公理与基本例子}
由整数、多项式、矩阵和函数的逐点运算引出环公理，回顾环的加法必须组成阿贝尔群。区分交换环与非交换环，检查子环是否包含原环的单位元。以 $\mathbb Z/n\mathbb Z$ 的剩余类运算说明同一集合可同时承载加法群和环结构。

\section{单位、零因子与整环}
定义单位（unit）、单位群 $R^\times$、左零因子和右零因子；在交换环中统一讨论零因子。定义整环（integral domain）为非零、交换且没有非零零因子的环，证明整环中的消去律。以矩阵环、整数环和模合数剩余类环区分可逆性与非零性。

\section{域、特征与素域}
定义域为非零且每个非零元素可逆的交换环。定义环的特征（characteristic）$\operatorname{char}R$ 为使 $n1_R=0$ 的最小正整数 $n$，不存在时记为 $0$。证明整环的特征为 $0$ 或素数，有限整环必为域。对素数 $p$ 记 $\mathbb F_p=\mathbb Z/p\mathbb Z$，并证明任意域的最小子域同构于 $\mathbb Q$ 或某个 $\mathbb F_p$。

\section{环同态与结构保持}
定义保持加法、乘法与单位元的环同态及环同构。比较自然剩余类映射、包含映射和多项式求值映射。说明环同态的核对加法形成子群，但一般不是本书约定下的含幺子环；这正是引入理想的动机。

\section{例题、练习与自检}
计算 $\mathbb Z/12\mathbb Z$ 的单位与非零零因子，验证高斯整数环 $\mathbb Z[i]=\{a+bi:a,b\in\mathbb Z\}$ 为整环。练习要求给出非交换环、非域整环及非整环交换环，并检查给定映射是否满足全部环同态条件。

